%% kapa: exam=1 round=27 year=2016 period=1 form=A subject=civil_law subject_ko=민법 seq=36 number=36
%% source: Q-Net 감정평가사 기출문제 (cst00309) artlSeq=5169402 file=extracted/1008/5169402/제27회 감정평가사 1차 문제/제27회 감정평가사 1교시문제(A형).pdf
%% conversion: kordoc 4.24.51 -> tools/gen1.py (auto); review=proofread-sonnet; flags=none
\begin{question}{36}
甲은 자신의 건물에 乙 명의의 전세권(전세금 1억원)을 설정해 주었다. 그 후 乙이 그 전세권에 丙 명의의 저당권(피담보채권액 7천만원)을 설정해 주었다. 이에 관한 설명으로 옳은 것을 모두 고른 것은? (다툼이 있으면 판례에 따름)
\begin{framedtext}
ㄱ. 乙의 전세권이 존속기간 만료로 종료된 경우, 그 전세권의 용익물권적 권능은 말소등기 없이도 당연히 소멸한다. \par
ㄴ. 乙의 전세권이 법정갱신되는 경우, 전세기간에 대한 변경등기 없이도 갱신된다. \par
ㄷ. 丙의 전세권저당권은 피담보채권을 수반하더라도 양도할 수 없다. \par
ㄹ. 乙의 전세권이 존속기간 만료로 종료된 경우, 丙은 전세권 자체에 대하여 저당권을 실행할 수 없게 된다. \par
\end{framedtext}
\begin{choices}
  \item ㄱ
  \item ㄴ, ㄷ
  \item ㄴ, ㄹ
  \item ㄱ, ㄴ, ㄹ
  \item ㄱ, ㄷ, ㄹ
\end{choices}
\end{question}
\answer{4}
